\documentclass[11pt]{article}
\usepackage[margin=1in]{geometry}
% Common macros for the anonymity series (shared across papers).
% Keep this file small: each paper should still compile as a standalone artifact.

\usepackage{amsmath,amssymb,amsthm}
\usepackage{mathtools}
\usepackage{microtype}
\usepackage{enumitem}
\usepackage{booktabs}
\usepackage{hyperref}
\usepackage[nameinlink,noabbrev]{cleveref}

\newcommand{\RR}{\mathbb{R}}
\newcommand{\NN}{\mathbb{N}}
\newcommand{\EE}{\mathbb{E}}
\newcommand{\PP}{\mathbb{P}}

\DeclareMathOperator{\TV}{TV}
\DeclareMathOperator{\KL}{D_{\mathrm{KL}}}
\DeclareMathOperator{\MI}{I}

\newcommand{\one}{\mathbf{1}}

\newcommand{\defeq}{\stackrel{\mathrm{def}}{=}}

\newtheorem{theorem}{Theorem}
\newtheorem{proposition}{Proposition}
\newtheorem{lemma}{Lemma}
\newtheorem{corollary}{Corollary}
\theoremstyle{definition}
\newtheorem{definition}{Definition}
\theoremstyle{remark}
\newtheorem{remark}{Remark}

% A small "callout" macro for short papers.
\newcommand{\takeaway}[1]{\paragraph{Takeaway.}\emph{#1}}


\title{Odds-Inflation Anonymity:\\ Posterior Odds Budgets, Composition, and Operational Limits}
\author{}
\date{March 21, 2026 (archive v499)}

\begin{document}
\maketitle

\begin{abstract}
We formalize anonymity as a bound on \emph{posterior odds inflation}: after observing a transcript, how many bits can the odds in favor of any single identity increase relative to a prior?
This pointwise notion yields referee-checkable, per-output guarantees that (i) are closed under post-processing, (ii) compose sequentially under interaction, and (iii) imply concrete limits on identification, shortlisting, and property inference.
We relate odds-inflation anonymity to pointwise maximal leakage and maximal leakage, and give practical conversions from common ``equalization'' contracts (e.g.\ total-variation closeness to a background distribution and R\'enyi-moment bounds).
The result is a compact endpoint metric that cleanly consumes mechanism-specific toolboxes (e.g.\ metadata-hiding overlay compilers) while stopping \emph{before} release packaging, repeated-use horizon accounting, and deployment-specific certificate serialization, which are owned elsewhere in the series.
\end{abstract}
\paragraph{Scope.}
We define a compact, mechanism-agnostic endpoint for anonymity: a bound on realized posterior odds inflation for every observable transcript.
The goal is a referee-checkable per-output statement that composes under interaction and can be consumed by overlay-specific ``compiler'' analyses.

\paragraph{Units.} Throughout, information budgets are measured in bits: we write $\log_2$ and $2^{\epsilon}$ for realized odds-inflation / pointwise-maximal-leakage caps. Imported results stated in nats can be converted by dividing by $\ln 2$.

\paragraph{Imports.}
Mechanism-specific contracts (scheduling/compilers, random-walk delegation, and stop-time padding) are imported from the published backbone notes and cited only as black-box prerequisites.
In particular, we treat termination-time hiding and TV certificates as coming from the scheduling, spectral, stop-time, and knapsack/transport papers.\cite{series:schedcompiler,series:spectral,series:stoptimeaddendum,series:knapsacktransport}
To keep the archive non-redundant, we also adopt the series contract-object interface (Synthesis~0): this paper exports an odds-inflation endpoint label, while other notes export schedule equalization, threat windows, and auditing objects that we consume via conversions.\cite{series:contractobjects}
We treat the bounded transcript projection and the repetition horizon as explicit claim fields: \texttt{exposure\_nf\_id} (ENF-ID) and \texttt{tw\_id} (TW-ID), imported from the trace-interface and threat-window synthesis notes.\cite{series:traceifaces,series:threatwindows}
Notation conventions for keys, state, transcripts, and transcript projections follow the series notation sheet (Synthesis~8).\cite{series:notation}


\paragraph{Exports.}
This paper exports (i) odds-inflation anonymity (Def.~\ref{def:odds_inflation}) as an endpoint metric, (ii) sequential/interactive composition bounds, and (iii) tight conversions from common equalization contracts (TV and moment bounds) into odds-inflation budgets.
These outputs are designed to be ``plugged into'' Boss Fight accounting and certificate workflows elsewhere in the series.\cite{series:endpointbridge}

\paragraph{Boundary.}
This note stops at the \emph{endpoint semantics}: which transcript projection, prior/context declaration, and per-output budget are being claimed, plus how common witness styles convert into that budget.
It does \emph{not} own long-horizon repeated-use sizing (Boss Fight~A), anonymous-DHT profiling specializations (Anonymity~B), or release-facing claim packaging and lineage objects (Anonymity~C and Release~A).\cite{series:bossfightA,series:profiling,series:abomoinl,series:releaseA}

\paragraph{Later-terse-paper citation posture.}
Later terse papers should cite this note only when the live issue is the endpoint anonymity contract itself: which declared transcript projection and threat window are in scope, which prior/context declaration governs the claim, which odds-inflation envelope (for example $\varepsilon_\infty$ or $(\varepsilon,\delta)$) is being published, or which generic TV / moment conversion produced that endpoint number.
If the live issue is instead repeated-use horizon sizing, stop at Boss Fight~A; if it is the anonymous-DHT witness specialization, stop at Anonymity~B; and if it is notarized claim packaging, lineage, or public release objects, stop at Anonymity~C and Release~A rather than widening this endpoint note into a deployment receipt.\cite{series:bossfightA,series:profiling,series:abomoinl,series:releaseA}

\paragraph{Artifact policy.}
This archive is source-only; supporting artifacts and tooling are available on request.\cite{series:artifactpolicy}


\section{Introduction}
Anonymous systems (including anonymous DHT lookups) often come with a \emph{stack} of mechanism-specific arguments:
timing padding, cover traffic, routing randomization, batching, and other defenses.
In our series, these are treated as modular ``compiler'' components for metadata-hiding overlays.
This paper isolates a short, mechanism-agnostic endpoint: an upper bound on the \emph{posterior odds} assigned to any identity after observing whatever transcript the system exposes.

Classical information-theoretic anonymity measures (e.g.\ entropy-based metrics) are valuable for intuition, but for certification we want a per-output bound that composes cleanly under interaction; odds inflation gives such an endpoint \cite{SerjantovDanezis2002}.

\paragraph{Non-redundancy with the series.}
We intentionally avoid re-deriving overlay- or routing-specific contracts here.
Instead, we prove generic theorems showing how \emph{any} component that exports a suitable privacy contract (likelihood-ratio caps, TV equalization, or moment bounds) can be converted into an odds-inflation guarantee that composes under interaction.
Mechanism details (stop-time padding, scheduling compilers, spectral delegation bounds) are cited as prerequisites, while repeated-use horizon sizing and release packaging are deferred to the Boss Fight and release-facing companions.

\section{Model and definitions}
Let $I\in\{1,\dots,N\}$ be a hidden identity with prior $\pi$ and let $O$ be any public observation (single output or full transcript).
Write $P_i$ for the law of $O$ given $I=i$, and $P=\sum_i \pi(i)P_i$ for the mixture (the public marginal).
We use $\TV(P,Q)=\tfrac12\|P-Q\|_1$ for total variation on finite alphabets.

\begin{definition}[Per-identity and worst-case odds inflation]\label{def:odds_inflation}
For any $o$ with $P(o)>0$, define the \emph{per-identity odds inflation}
\[
L_i(o)\ \defeq\ \log_2\frac{\pi(i\mid o)}{\pi(i)}\ =\ \log_2\frac{P_i(o)}{P(o)}.
\]
Define the \emph{realized (worst-case) odds inflation}
\[
\ell(o)\ \defeq\ \max_{i\in\{1,\dots,N\}} L_i(o).
\]
\end{definition}

\begin{definition}[$\varepsilon_\infty$ and $(\varepsilon,\delta)$ odds-inflation anonymity]\label{def:epsinf}
A mechanism is \emph{$\varepsilon_\infty$-anonymous} (relative to prior $\pi$) if $\ell(o)\le \varepsilon$ for all $o$ with $P(o)>0$.
More generally, it is \emph{$(\varepsilon,\delta)$-anonymous} if for every event $A$,
\begin{equation}\label{eq:hockeystick}
P_i(A)\ \le\ 2^\varepsilon P(A) + \delta
\qquad\text{for all } i\in\{1,\dots,N\}.
\end{equation}
\end{definition}

\begin{remark}[Interpretation]
$\ell(o)$ is the number of \emph{bits} by which the observation $o$ can multiply the posterior odds of the most-favored identity (base 2). Results phrased in nats use $\ln$ and can be converted by dividing by $\ln 2$.
The $(\varepsilon,\delta)$ form \eqref{eq:hockeystick} is a tail envelope: with probability at most $\delta$ (under $P_i$), the realized loss can exceed $\varepsilon$ by an arbitrary amount.
Equivalently, \eqref{eq:hockeystick} is the hockey-stick (\emph{$E_\gamma$}) divergence bound $E_{2^\varepsilon}(P_i\|P)\le\delta$, i.e., $\sup_A\{P_i(A)-2^\varepsilon P(A)\}\le\delta$ \cite{LiuCuffVerdu2017}.
Hockey-stick divergence also underlies privacy-profile style accounting in differential privacy and remains an active object of study for auditing and restricted-adversary models (including ``measured'' hockey-stick divergence in quantum settings) \cite{NuradhaMeasuredHockeyStick2025}.
This mirrors the usual $(\varepsilon,\delta)$ differential privacy inequality but with a mixture baseline \cite{DworkRoth2014}.
Related inferential-privacy frameworks (e.g., Pufferfish) are also phrased as posterior-odds ratio bounds between specified secret pairs; odds-inflation anonymity can be viewed as a one-vs-mixture specialization tailored to anonymity/attribution settings \cite{KiferMachanavajjhala2014Pufferfish}.
\end{remark}

\section{Priors and contexts as certificate inputs}\label{sec:priors}
Odds inflation is defined relative to a prior $\pi$ because the public marginal $P=\sum_i \pi(i)P_i$ appears in the Bayes factor $P_i/P$.
This is a feature: a deployment's \emph{context} (who could plausibly be the requester) is part of the privacy surface.
Accordingly, in our series we treat priors and eligibility constraints as first-class certificate inputs (Anonymity~C) \cite{series:abomoinl}.

\paragraph{Certification practice.}
A paper should state which prior or prior family its claimed budget targets.
Two common, conservative patterns are:
\begin{enumerate}[leftmargin=*]
\item \textbf{Uniform over an eligibility pool.} Declare an eligibility interface (e.g.\ ``clients active in this epoch''), and take $\pi$ uniform over that pool.
Then the operational corollaries in \Cref{prop:shortlist,prop:property} translate directly into lower bounds on the number of plausible requesters.
\item \textbf{Capped prior mass.} Declare a cap $\pi(i)\le \pi_{\max}$ for all identities in scope.
Under $\varepsilon_\infty$-anonymity, every posterior mass satisfies $\pi(i\mid o)\le 2^{\varepsilon}\pi_{\max}$, so deployments can report a policy-friendly statement such as ``no single identity exceeds $p^\star$ posterior mass'' by enforcing $2^{\varepsilon}\pi_{\max}\le p^\star$.
\end{enumerate}
In both cases, \emph{eligibility metadata must be treated as a public interface} and budgeted (or hidden) rather than silently absorbed into the prior.

\begin{proposition}[Eligibility/selection disclosures add a prior-shift term]\label{prop:selection}
Let $C\subseteq\{1,\dots,N\}$ with $\pi(C)>0$ and let $E\defeq\one\{I\in C\}$ denote an eligibility/selection bit that becomes public (e.g.\ an observer learns which shard/committee the requester must belong to).
Let $P_C(o)\defeq\sum_{j\in C}\pi(j\mid C)P_j(o)$ be the public marginal conditioned on $I\in C$.
Then for any $i\in C$ and any observation $o$ with $P_C(o)>0$, the realized odds inflation for the combined disclosure $(E{=}1,O{=}o)$ satisfies
\[
\log_2\frac{\pi(i\mid E{=}1,o)}{\pi(i)}\ =\ -\log_2 \pi(C)\ +\ \log_2\frac{P_i(o)}{P_C(o)}.
\]
In particular, if the mechanism is $\varepsilon_\infty$-anonymous relative to the conditional prior $\pi(\cdot\mid C)$ (i.e., $\max_{i\in C}\log_2\tfrac{P_i(o)}{P_C(o)}\le\varepsilon$ for all $o$), then relative to the original prior it is $(\varepsilon+\log_2\tfrac{1}{\pi(C)})_\infty$-anonymous for identities in $C$ (equivalently, $\varepsilon-\log_2\pi(C)$).
Under a uniform prior this adds the familiar term $\log_2(N/|C|)$.
\end{proposition}

\begin{proof}
For $i\in C$, $P_i(E{=}1)=1$ while $P(E{=}1)=\pi(C)$.
Moreover $P(E{=}1,O=o)=\pi(C)\,P_C(o)$.
Therefore
$\log_2\tfrac{P_i(E{=}1,O=o)}{P(E{=}1,O=o)}=\log_2\tfrac{P_i(o)}{\pi(C)P_C(o)}=-\log_2\pi(C)+\log_2\tfrac{P_i(o)}{P_C(o)}$.
Finally, $\pi(i\mid E{=}1,o)/\pi(i)=P_i(E{=}1,O=o)/P(E{=}1,O=o)$ by Bayes.
\end{proof}

\begin{remark}[Context-aware leakage bounds]
Recent pointwise-maximal-leakage analyses show that even for fixed mechanisms, leakage can be substantially smaller under a committed prior model (e.g., minimum-probability assumptions) than under worst-case, context-free bounds \cite{ZhaoSaeidianOechtering2026ContextAware}.
Odds-inflation anonymity makes this dependence explicit by treating the prior (or prior family) as a certificate input.
\end{remark}

\section{Relation to PML/MaxL (pointer)}
The realized odds inflation $\ell(o)=\max_i\log_2\tfrac{P_i(o)}{P(o)}$ is exactly pointwise maximal leakage (PML) at outcome $o$ under the same prior; maximal leakage (MaxL) is the log-moment of $\exp(\ell(O))$ and yields the standard $\log_2(\alpha K)$ sizing rule for exact identification.
To keep this paper non-redundant, we treat these equivalences as ``owned'' by the endpoint-bridge synthesis note and the Boss-Fight budgets paper; see \cite{series:endpointbridge,series:bossfightA}.

\section{Closure properties}
\subsection{Post-processing}
\begin{theorem}[Post-processing]\label{thm:postprocessing}
Let $Z=g(O)$ be any (possibly randomized) function of $O$.
Then for every realized value $z$,
\[
\ell_Z(z)\le \operatorname*{ess\,sup}_{o:\,K(z\mid o)>0}\,\ell_O(o),
\]
where $K$ is the Markov kernel of $g$ (for deterministic $g$, this is $\sup_{o:\,g(o)=z}\ell_O(o)$).
In particular, $\varepsilon_\infty$-anonymity and $(\varepsilon,\delta)$-anonymity are preserved under post-processing.
\end{theorem}

\begin{proof}
Fix $z$ with $P_Z(z)>0$. For each $i$, write the Bayes factor $B_i(o)\defeq P_i(o)/P(o)$.
Then
\[
\frac{P_i(Z=z)}{P(Z=z)}=\frac{\sum_o P(o)B_i(o)\,K(z\mid o)}{\sum_o P(o)\,K(z\mid o)}=\EE_P\!\bigl[B_i(O)\mid Z=z\bigr],
\]
where $K$ is the Markov kernel of $g$ (randomized post-processing) and the conditional expectation is under the mixture $P$.
Therefore
\[
\max_i \frac{P_i(Z=z)}{P(Z=z)} \le \EE_P\!\Bigl[\max_i B_i(O)\,\Bigm|\, Z=z\Bigr].
\]
Since a conditional expectation is bounded by the essential supremum on the conditioning event,
\[
\EE_P\!\Bigl[\max_i B_i(O)\,\Bigm|\, Z=z\Bigr]
\le \operatorname*{ess\,sup}_{o:\,K(z\mid o)>0}\,\max_i B_i(o)
=\exp\Bigl(\operatorname*{ess\,sup}_{o:\,K(z\mid o)>0}\,\ell_O(o)\Bigr).
\]
Taking logs yields the stated bound.
For $(\varepsilon,\delta)$, note that \eqref{eq:hockeystick} is the $E_{2^\varepsilon}$ (hockey-stick) divergence bound $E_{2^\varepsilon}(P_i\|P)\le\delta$, and $E_\gamma$ satisfies data processing under Markov kernels, so the same $(\varepsilon,\delta)$ pair holds for $Z$ \cite{LiuCuffVerdu2017}.
\end{proof}

\subsection{Sequential composition under interaction}
Consider an interactive transcript $O_{1:T}$ where each step may depend on the past.
Write $\ell_t$ for the realized odds inflation after observing prefix $O_{1:t}$.

\begin{theorem}[Sequential composition for $\varepsilon_\infty$]\label{thm:epsinf_comp}
If each step $t$ satisfies the conditional per-step bound
\[
\max_i \log_2\frac{P_i(O_t\mid O_{1:t-1})}{P(O_t\mid O_{1:t-1})}\ \le\ \varepsilon_t
\quad \text{almost surely},
\]
then the full transcript is $\bigl(\sum_{t=1}^T \varepsilon_t\bigr)_\infty$-anonymous.
\end{theorem}

\begin{proof}
Let $B_{i,t}\defeq \frac{P_i(O_{1:t})}{P(O_{1:t})}$ be the Bayes factor for identity $i$ after $t$ steps.
By the chain rule,
\[
B_{i,T}=\prod_{t=1}^T \frac{P_i(O_t\mid O_{1:t-1})}{P(O_t\mid O_{1:t-1})}.
\]
Taking $\log_2$, maximizing over $i$, and applying the per-step bound yields
$\ell(O_{1:T})=\log_2\max_i B_{i,T}\le \sum_{t=1}^T \varepsilon_t$ almost surely.
\end{proof}

\begin{theorem}[Sequential composition for $(\varepsilon,\delta)$ checkpoints]\label{thm:epsdelta_comp}
Consider an interactive transcript $O_{1:T}$.
Assume that for each step $t$ there exist numbers $(\varepsilon_t,\delta_t)$ such that, almost surely over the public history $O_{1:t-1}$,
for every event $A$ and every identity $i$,
\begin{equation}\label{eq:step_epsdelta}
P_i(O_t\in A\mid O_{1:t-1})\ \le\ 2^{\varepsilon_t}\,P(O_t\in A\mid O_{1:t-1})\ +\ \delta_t,
\end{equation}
where $P(\cdot\mid O_{1:t-1})$ is the mixture (posterior) baseline at step $t$.
Then the full transcript satisfies, for every event $B$ and every identity $i$,
\[
P_i(O_{1:T}\in B)\ \le\ 2^{\varepsilon}\,P(O_{1:T}\in B)\ +\ \delta_{\mathrm{tot}},
\qquad \varepsilon:=\sum_{t=1}^T \varepsilon_t,\quad
\delta_{\mathrm{tot}}:=\sum_{t=1}^T \delta_t\,\exp\Big(\sum_{s<t}\varepsilon_s\Big).
\]
In particular, if $\varepsilon_t\ge 0$ for all $t$, then $\delta_{\mathrm{tot}}\le 2^{\varepsilon}\sum_t\delta_t$.
\end{theorem}

\begin{proof}
We prove the bound by induction on $T$.
For $T=1$ it is \eqref{eq:step_epsdelta}.
Assume it holds up to $T-1$.
Fix an event $B$ on $O_{1:T}$ and write its sections
$B_h:=\{o_T:\ (h,o_T)\in B\}$ for histories $h=o_{1:T-1}$.
By the tower property,
\[
P_i(B)=\EE_{P_i}\bigl[P_i(O_T\in B_{O_{1:T-1}}\mid O_{1:T-1})\bigr].
\]
Applying \eqref{eq:step_epsdelta} with $A=B_{O_{1:T-1}}$ yields
\[
P_i(B)\le 2^{\varepsilon_T}\,\EE_{P_i}[\phi(O_{1:T-1})]+\delta_T,
\qquad \phi(h):=P(O_T\in B_h\mid O_{1:T-1}=h).
\]
To relate $\EE_{P_i}[\phi(O_{1:T-1})]$ to $\EE_P[\phi(O_{1:T-1})]=P(B)$, use the layer-cake identity
$\phi=\int_0^1 \one\{\phi\ge u\}\,du$ to reduce the claim to events and apply the induction hypothesis to each threshold set $\{\phi\ge u\}$.
This yields
\[
\EE_{P_i}[\phi(O_{1:T-1})]\le 2^{\sum_{t<T}\varepsilon_t}\,\EE_P[\phi(O_{1:T-1})]+\sum_{t<T}\delta_t\,\exp\Big(\sum_{s<t}\varepsilon_s\Big).
\]
Combining and simplifying yields the stated $\delta_{\mathrm{tot}}$.
\end{proof}

\begin{corollary}[Homogeneous checkpoints]\label{cor:homog_checkpoints}
If the per-step bound \eqref{eq:step_epsdelta} holds with the same pair $(\varepsilon_t,\delta_t)\equiv(\varepsilon,\delta)$ for all $t$ (uniformly over histories), then for every event $B$ and every identity $i$,
\[
P_i(O_{1:T}\in B)\ \le\ 2^{T\varepsilon}\,P(O_{1:T}\in B)\ +\ \delta\,\sum_{j=0}^{T-1} 2^{j\varepsilon}
\ =\ 2^{T\varepsilon}P(B)+\begin{cases}
\delta\,\dfrac{2^{T\varepsilon}-1}{2^{\varepsilon}-1}, & \varepsilon\neq 0,\\
T\delta, & \varepsilon=0.
\end{cases}
\]
In particular, if $\varepsilon\ge 0$ then the slack is at most $2^{T\varepsilon}\,T\delta$.
\end{corollary}

\begin{remark}[Anytime-valid monitoring]
The per-identity Bayes factor process $B_{i,t}=\exp(L_i(O_{1:t}))$ is a nonnegative martingale under the mixture law.
This connects odds inflation to \emph{e-processes} and time-uniform bounds (Ville-style inequalities), and to privacy odometers/filters in adaptive DP composition.
See \cite{HowardRamdas2018,RogersOdometers2016}.
\end{remark}

\section{Operational consequences}
Odds caps imply concrete limits on what any adversary can do with the observation.

\begin{lemma}[Suspect sets must have prior mass]\label{lem:suspect}
For any threshold $\tau>0$ and any observation $o$,
\[
\pi\bigl(\{i:\ L_i(o)\ge \log_2\tau\}\bigr)\ \le\ \frac{1}{\tau}.
\]
In particular, if $\ell(o)\le \varepsilon$ then every identity $i$ satisfies $\pi(i\mid o)\le 2^\varepsilon \pi(i)$.
\end{lemma}

\begin{proposition}[Shortlists and credible sets]\label{prop:shortlist}
Suppose the mechanism is $\varepsilon_\infty$-anonymous.
Let $S(o)\subseteq\{1,\dots,N\}$ be any post-processed shortlist output on observation $o$.
If $S(o)$ is required to have posterior coverage $\pi(S(o)\mid o)\ge 1-\beta$ for all $o$, then $\pi(S(o))\ge 2^{-\varepsilon}(1-\beta)$ for all $o$.
Under a uniform prior this forces $|S(o)|\ge (1-\beta)2^{-\varepsilon}N$.
\end{proposition}

\begin{proof}
Fix any observation $o$ with $P(o)>0$.
By $\varepsilon_\infty$-anonymity, for every identity $i$ we have $\pi(i\mid o)\le 2^{\varepsilon}\pi(i)$.
Summing over $i\in S(o)$ gives
\[
\pi(S(o)\mid o)=\sum_{i\in S(o)}\pi(i\mid o)\ \le\ 2^{\varepsilon}\sum_{i\in S(o)}\pi(i)=2^{\varepsilon}\pi(S(o)).
\]
If $\pi(S(o)\mid o)\ge 1-\beta$, then rearranging yields $\pi(S(o))\ge 2^{-\varepsilon}(1-\beta)$.
For the uniform prior, $\pi(S(o))=|S(o)|/N$, hence $|S(o)|\ge (1-\beta)2^{-\varepsilon}N$.
\end{proof}

\begin{corollary}[Exact identification is bounded]\label{cor:identification}
Assume a uniform prior on $\{1,\dots,N\}$ and $\varepsilon_\infty$-anonymity.
Then for any estimator $\widehat I=\widehat I(O)$,
\[
\PP(\widehat I=I)\ \le\ \frac{2^{\varepsilon}}{N}.
\]
\end{corollary}

\begin{proof}
Condition on $O=o$.
The best success probability is the MAP mass $\max_i\pi(i\mid o)$, which is at most $2^{\varepsilon}/N$ under $\varepsilon_\infty$-anonymity and the uniform prior.
Average over $o$.
\end{proof}


\begin{corollary}[Identification under $(\varepsilon,\delta)$]\label{cor:identification_epsdelta}
Assume a uniform prior on $\{1,\dots,N\}$ and $(\varepsilon,\delta)$-anonymity in the sense of \Cref{def:epsinf}.
Then for any estimator $\widehat I=\widehat I(O)$,
\[
\PP(\widehat I=I)\ \le\ \frac{2^{\varepsilon}}{N}+\delta.
\]
\end{corollary}

\begin{proof}
Let $A_i=\{o:\ \widehat I(o)=i\}$.
Under the uniform prior, $\PP(\widehat I=I)=\frac{1}{N}\sum_{i=1}^N P_i(A_i)$.
Apply \eqref{eq:hockeystick} to each $A_i$ to obtain $P_i(A_i)\le 2^\varepsilon P(A_i)+\delta$.
Averaging over $i$ and using $\sum_i P(A_i)=1$ yields the claim.
\end{proof}



\begin{proposition}[Property inference is also capped]\label{prop:property}
Assume $\varepsilon_\infty$-anonymity. For any subset $E\subseteq\{1,\dots,N\}$ and any observation $o$ with $P(o)>0$,
\[
\pi(E\mid o)\ \le\ 2^{\varepsilon}\,\pi(E).
\]
Equivalently, no post-processed binary property of the identity can have its posterior probability inflated by more than $2^{\varepsilon}$.
\end{proposition}

\begin{proof}
Sum the pointwise bounds $\pi(i\mid o)\le 2^{\varepsilon}\pi(i)$ over $i\in E$.
\end{proof}

\section{Practical conversions from common contracts}
This section records two conversion lemmas that let mechanism-specific analyses ``plug into'' odds inflation.
The TV-equalization form is the same interface used by our TV-privacy padding compilers and dual-certificate tooling \cite{series:knapsacktransport}.

\subsection{TV equalization to a smoothed reference}
Let $Q$ be a public reference distribution on outcomes with $Q_{\min}:=\min_o Q(o)>0$.
In many anonymity channels, a mechanism exports a \emph{common-reference} guarantee: both the per-identity output law $P_i$ and the public marginal $P=\sum_i \pi(i)P_i$ are close to the same $Q$ (e.g.\ via equalization, smoothing, or explicit cover traffic).
This suffices to bound the identity-vs-mixture Bayes factors $P_i/P$ that define odds inflation.

\begin{lemma}[Common-reference TV-to-likelihood-ratio]\label{lem:tv_to_lr}
Let $P$ and $R$ be distributions on a finite alphabet, and let $Q$ satisfy $Q_{\min}>0$.
If $\TV(P,Q)\le \delta$ and $\TV(R,Q)\le \delta$ with $\delta<Q_{\min}/2$, then for all outcomes $o$,
\[
\left|\log_2\frac{P(o)}{R(o)}\right|\ \le\ \log_2\frac{Q_{\min}+2\delta}{Q_{\min}-2\delta}.
\]
\end{lemma}

\begin{proof}
Since $\|P-Q\|_1\le 2\delta$, for each $o$ we have $P(o)\le Q(o)+2\delta$.
Similarly $\|R-Q\|_1\le 2\delta$ implies $R(o)\ge Q(o)-2\delta$.
Therefore
\[
\frac{P(o)}{R(o)}\ \le\ \frac{Q(o)+2\delta}{Q(o)-2\delta}
\ \le\ \frac{Q_{\min}+2\delta}{Q_{\min}-2\delta},
\]
where the last step uses that $(x+2\delta)/(x-2\delta)$ is decreasing for $x>2\delta$.
The reverse ratio $R(o)/P(o)$ is handled symmetrically.
\end{proof}

\begin{corollary}[Equalization $\Rightarrow$ per-output odds-inflation cap]\label{cor:eq_to_epsinf}
If a mechanism satisfies $\TV(P_i,Q)\le \delta$ for all identities $i$ (with the same $Q$ and $\delta$),
then $\TV(P,Q)\le \delta$ for the public marginal $P=\sum_i \pi(i)P_i$.
Consequently, if $\delta<Q_{\min}/2$ then for every realized outcome $o$,
\[
\ell(o)=\max_i \log_2\frac{P_i(o)}{P(o)}\ \le\ \log_2\frac{Q_{\min}+2\delta}{Q_{\min}-2\delta}.
\]
\end{corollary}

\begin{proof}
By convexity of TV, $\TV(P,Q)=\TV(\sum_i\pi(i)P_i,Q)\le \sum_i \pi(i)\TV(P_i,Q)\le \delta$.
Apply \Cref{lem:tv_to_lr} with $P=P_i$ and $R=P$ and maximize over $i$.
\end{proof}

\begin{remark}
This conversion is conservative but operationally important: without a positive background mass ($Q_{\min}=0$), worst-case odds inflation can be infinite.
For anonymous DHTs, this is the ``rare committee'' failure mode.
When $Q_{\min}$ is too small to enforce, use a trimmed-support tail guarantee instead (next lemma).
\end{remark}

\begin{lemma}[Common-reference trimmed-support envelope]\label{lem:tv_trimmed}
Let $\tau>0$ and define the ``well-supported'' set $S_\tau\defeq\{o:\ Q(o)\ge \tau\}$.
If $\TV(P,Q)\le \delta$ and $\TV(R,Q)\le \delta$ with $\delta<\tau/2$, then for all $o\in S_\tau$,
\[
\left|\log_2\frac{P(o)}{R(o)}\right|\ \le\ \log_2\frac{\tau+2\delta}{\tau-2\delta}.
\]
Moreover, $P(S_\tau^c)\le Q(S_\tau^c)+\delta$ and $R(S_\tau^c)\le Q(S_\tau^c)+\delta$.
Thus a TV equalization contract yields a finite loss bound on $S_\tau$ and pushes the remaining probability mass into explicit ``rare outcome'' tail terms.
\end{lemma}

\subsection{Moment (R\'enyi) bounds to tail envelopes}
Mechanism analyses often bound conditional R\'enyi divergences or log-moments of the privacy loss.
Such bounds convert directly into $(\varepsilon,\delta)$ odds-inflation envelopes by a one-line Chernoff argument.

\begin{lemma}[R\'enyi-to-tail conversion]\label{lem:renyi_tail}
Fix an identity $i$ and let $L_i(o)=\log_2\frac{P_i(o)}{P(o)}$ be the per-output privacy loss.
If for some order $\alpha>1$,
\[
D_\alpha(P_i\|P)\ \le\ \rho_\alpha,
\]
then for every threshold $\varepsilon$,
\[
P_i\!\bigl(L_i(O)\ge \varepsilon\bigr)\ \le\ \exp\bigl((\alpha-1)(\rho_\alpha-\varepsilon)\bigr).
\]
Equivalently, with $\delta=\exp((\alpha-1)(\rho_\alpha-\varepsilon))$, the odds inflation exceeds $\varepsilon$ with probability at most $\delta$ under $P_i$.
\end{lemma}

\begin{proof}
By definition of R\'enyi divergence and the identity $\exp((\alpha-1)L_i)=\bigl(P_i/P\bigr)^{\alpha-1}$,
\[
\EE_{P_i}\!\left[\exp\!\bigl((\alpha-1)L_i(O)\bigr)\right]
=\sum_o P_i(o)\left(\frac{P_i(o)}{P(o)}\right)^{\alpha-1}
=\sum_o P_i(o)^{\alpha}\,P(o)^{1-\alpha}
=\exp\!\bigl((\alpha-1)D_\alpha(P_i\|P)\bigr)
\le 2^{(\alpha-1)\rho_\alpha}.
\]
Applying Markov's inequality under $P_i$ gives
\[
P_i(L_i\ge \varepsilon)=P_i\!\left(\exp((\alpha-1)L_i)\ge 2^{(\alpha-1)\varepsilon}\right)
\le 2^{(\alpha-1)(\rho_\alpha-\varepsilon)}.
\]
\end{proof}

\begin{remark}[Connection to RDP/moment accountants]
Lemma~\ref{lem:renyi_tail} is the same tail conversion that underlies R\'enyi differential privacy (RDP) and moment-accountant style envelopes, but applied to the identity-vs-mixture loss $L_i$ rather than neighboring datasets.
See \cite{MironovRDP2017} for the RDP framing, \cite{DworkRoth2014} for differential privacy background, and \cite{RogersOdometers2016} for adaptive composition tooling.
\end{remark}

\section{A worked audit-style example (telemetry to odds inflation)}\label{sec:worked_example}
Suppose a deployment fixes a discrete interface (output alphabet of size $K$) and a public background $Q$ with floor $Q_{\min}$ (or a trimmed floor $\tau$).
For each audited identity/class $i$ (or each equivalence class of identities), an auditor collects a window of $m$ i.i.d.\ samples from the interface and forms an empirical distribution $\widehat P_i$.
A conservative $L_1$ deviation bound of Weissman et al.\ yields a confidence radius for $\|\widehat P_i-P_i\|_1$ from finite samples \cite{Weissman2003L1}.
Translating to TV gives an upper confidence bound on $\TV(P_i,Q)$ when $Q$ is known.
Once all audited $P_i$ satisfy $\TV(P_i,Q)\le\delta$, \Cref{cor:eq_to_epsinf} converts this equalization certificate into an odds-inflation cap against the public marginal.

\paragraph{Beyond TV radii (optional).}
If a deployment can support black-box hypothesis testing between $P_i$ and the mixture $P$, one can audit richer tradeoff curves (e.g., $f$-DP / tradeoff functions) rather than a single TV radius; see Askin et al.\ for recent black-box $f$-DP estimation and auditing methods \cite{AskinUSENIXSecurity2025}.

For a fully worked audit walk that threads telemetry, equalization, and odds-inflation semantics through the maintained example bundle, see Synthesis~17.\cite{series:workedexample}

Release-facing packaging of the resulting claim surface belongs to Anonymity~C and Release~A rather than to this note.\cite{series:abomoinl,series:releaseA}

This ``telemetry $\rightarrow$ equalization $\rightarrow$ odds inflation'' pipeline is what Anonymity~B specializes to anonymous DHT committee traces, and what Anonymity~C records in an ABOM claim object \cite{series:profiling,series:abomoinl}.

\section{How to consume this paper in the series}
Once a module exports either (i) a likelihood-ratio cap, (ii) TV equalization to a smoothed reference, or (iii) a moment bound,
the odds-inflation endpoint and its composition rules apply verbatim.
This allows each mechanism paper to avoid re-proving downstream operational consequences (shortlisting limits, identification bounds, etc.) and simply cite this paper.
For timing/termination-time contracts in anonymous overlay lookups, see the companion mechanism toolboxes \cite{series:schedcompiler,series:prefixcapacities}; for delegation-style overlays, see \cite{series:spectral}.


\paragraph{A minimal public odds-inflation endpoint card.}
Table~\ref{tab:minimal-odds-inflation-card} is the smallest public packet later terse papers should need when the live question is simply \emph{which endpoint anonymity contract was declared for this observation surface?}

\begin{table}[t]
\centering
\small
\begin{tabular}{@{}p{0.33\linewidth}p{0.57\linewidth}@{}}
\toprule
\textbf{Public row} & \textbf{Pinned value}\\
\midrule
Declared observation surface & \texttt{exposure\_nf\_id = enf-committee-contact-bucket-v2}; \texttt{tw\_id = tw-24h-200-lookups} \\
Prior / context declaration & uniform over $N=4096$ eligible requesters in the declared shard/context \\
Endpoint contract form & $\varepsilon_\infty = 0.75$ bits of posterior-odds inflation on the declared surface \\
Pointwise posterior cap & for every identity $i$ and every realized observation $o$, $\pi(i\mid o) \le 2^{0.75}\pi(i) \approx 1.682\,\pi(i)$ \\
Exact-identification cap & any estimator obeys $\Pr(\widehat I = I) \le 2^{0.75}/4096 \approx 4.11\times 10^{-4}$ \\
95\% shortlist floor & any shortlist with posterior coverage at least $0.95$ must contain at least $0.95\cdot 2^{-0.75}\cdot 4096 \approx 2314$ identities \\
Downstream owner & Boss Fight~A / Anonymity~B / ABOM--OINL / Release~A for horizon accounting, anonymous-DHT specialization, and lineage packaging \\
\bottomrule
\end{tabular}
\caption{A minimal public odds-inflation endpoint card. This is the non-decorative public answer later terse papers can quote directly when the live issue is the declared endpoint anonymity contract rather than repeated-use sizing, anonymous-DHT witness specialization, or release lineage.}
\label{tab:minimal-odds-inflation-card}
\end{table}

This card is intentionally narrow: it pins the declared observation surface, the prior/context declaration, the endpoint budget representation, and two referee-facing operational corollaries, but it does not absorb horizon accounting, mechanism-specific witness dossiers, or release objects. Those remain owned by the downstream companions already cited above.\cite{series:bossfightA,series:profiling,series:abomoinl,series:releaseA}

\paragraph{One tiny endpoint drill.}
With the card values above,
\[
\Pr(\widehat I = I) \le \frac{2^{0.75}}{4096} \approx 4.11\times 10^{-4},
\]
by Corollary~\ref{cor:identification}. Likewise, Proposition~\ref{prop:shortlist} gives
\[
|S(o)| \ge 0.95\cdot 2^{-0.75}\cdot 4096 \approx 2314
\]
for every shortlist whose posterior coverage is at least $95\%$. A later terse paper can quote the card and this arithmetic without reopening the conversion lemmas or the release stack.

\paragraph{One tiny same-number cutover drill.}
Keep the same published budget number $\varepsilon_\infty=0.75$ bits, but quietly change the prior/context row from ``uniform over $N=4096$ eligible requesters in the declared shard'' to ``uniform over $N=1024$ eligible requesters in a narrower shard.'' Then the old endpoint card no longer licenses the same operational arithmetic, even though the budget number itself has not moved: the exact-identification cap becomes
\[
\frac{2^{0.75}}{1024} \approx 1.64\times 10^{-3},
\]
which is four times larger than the old $4.11\times 10^{-4}$ cap, and the $95\%$ shortlist floor drops to
\[
0.95\cdot 2^{-0.75}\cdot 1024 \approx 579.
\]
The same fail-closed rule applies if a later note reuses ``$0.75$ bits'' while changing ENF-ID or TW-ID: the budget number by itself is not the public claim object. A later terse paper may reuse the old arithmetic only while the declared observation surface, threat window, and prior/context declaration all stay fixed.

\paragraph{Minimal reopen ladder.}
Once the live question is no longer ``what exact endpoint anonymity card was declared here?'', the first reopen owner should be named explicitly rather than inferred. Reopen Boss Fight~A when the repeated-use horizon, mixed witness stack, or spending rule changes while the same endpoint semantics stay in force. Reopen Anonymity~B when the anonymous-DHT equalization witness, secret scope, or committee-trace specialization changes under the same endpoint metric. Reopen State~1 when state-surface drift changes which prior/context or exposed control-plane surface the endpoint card is about. Reopen Operational~A or Operational~B when retry-discipline, replay-plan, or runtime-conformance questions change the declared projection itself. Reopen Eval~1 when the missing issue is notarized evaluator evidence or anti-grinding attempt logging. Reopen Anonymity~C and Release~A only when the live question has widened to claim-object packaging, signed public receipts, or cross-release lineage rather than endpoint semantics. This note therefore exports one compact stop rule: quote the endpoint card only while the declared surface, threat window, prior/context, and budget representation are the same public object.\cite{series:bossfightA,series:profiling,series:state1,series:operationalA,series:operationalB,series:eval1,series:abomoinl,series:releaseA}

\section{Takeaways}
\begin{itemize}[leftmargin=*]
\item \textbf{Endpoint metric:} odds inflation is a posterior-odds cap for a transcript projection $X$ and a declared prior/context. It is the right ``final number'' to publish when the operational question is \emph{how well can an adversary decide?}
\item \textbf{Composable by default:} sequential use over epochs composes additively in log-odds (and via standard odometer/filter reasoning when stopping rules are adaptive). Mechanism papers in this archive should export per-epoch or per-step budgets and cite this paper for the endpoint/composition rule, while Boss Fight~A owns the repeated-use horizon interpretation.\cite{series:bossfightA}
\item \textbf{Operational conversions:} odds budgets induce (i) shortlist limits and (ii) identification sizing rules of the form $\log_2(\alpha K)$, enabling ``privacy-to-knobs'' compilation without re-deriving decision-theoretic consequences (see also Synthesis~5).\cite{series:endpointbridge}
\item \textbf{Claim surface, not release packaging:} the key inputs are the transcript projection, prior/context declaration, and budget representation. ABOM/OINL (Anonymity~C) is the canonical carrier for committing to those semantics in a source-only release; Release~A then packages lineage, anti-equivocation, and signed public release objects.\cite{series:abomoinl,series:releaseA}
\end{itemize}

\begin{thebibliography}{99}\small

\bibitem{series:contractobjects}
Anonymous.
\newblock Contract Objects for Anonymous-DHT Privacy Budgets and Claims.
\newblock Series synthesis note (Synthesis~0), 2026. (This archive.)


\bibitem{series:traceifaces}
Anonymous.
\newblock Trace Interfaces for Anonymous DHTs: Observation Tiers, Committee-Contact Traces, and Exposure Normal Forms.
\newblock Series synthesis note (Synthesis~3), 2026. (This archive.)

\bibitem{series:threatwindows}
Anonymous.
\newblock Threat Models and Repetition Windows for Anonymous DHTs: A Short Declaration Checklist for Referee-Grade Budget Claims.
\newblock Series synthesis note (Synthesis~4), 2026. (This archive.)

\bibitem{series:endpointbridge}
Anonymous.
\newblock Endpoint Metrics Bridge for Anonymous Overlays: Odds Inflation, Pointwise Maximal Leakage, and Maximal Leakage.
\newblock Series synthesis note, 2026. (This archive.)

\bibitem{SerjantovDanezis2002}
A.~Serjantov and G.~Danezis.
\newblock Towards an information theoretic metric for anonymity.
\newblock In \emph{Privacy Enhancing Technologies (PET)}, 2002.

\bibitem{RogersOdometers2016}
R.~Rogers, A.~Roth, J.~Ullman, and S.~Vadhan.
\newblock Privacy odometers and filters: pay-as-you-go composition.
\newblock In \emph{NeurIPS}, 2016. arXiv:1605.08294.

\bibitem{HowardRamdas2018}
S.~R.~Howard, A.~Ramdas, J.~McAuliffe, and J.~Sekhon.
\newblock Time-uniform Chernoff bounds via nonnegative supermartingales.
\newblock \emph{Probability Surveys}, 18:29--86, 2021. arXiv:1808.03204.

\bibitem{ZhaoSaeidianOechtering2026ContextAware}
H.~Zhao, S.~Saeidian, and T.~J.~Oechtering.
\newblock Context-aware privacy bounds for linear queries.
\newblock arXiv:2601.02855 [cs.IT], Jan.\ 2026. \url{https://arxiv.org/abs/2601.02855}.

\bibitem{AskinUSENIXSecurity2025}
\"{O}.~Askin, H.~Dette, M.~Dunsche, T.~Kutta, Y.~Lu, Y.~Wei, and V.~Zikas.
\newblock General-purpose $f$-DP estimation and auditing in a black-box setting.
\newblock In \emph{USENIX Security}, 2025. Preprint: arXiv:2502.07066. \url{https://arxiv.org/abs/2502.07066}.

% Common bibliography items for already-published papers in the anonymity series.
% Keep these stable so other papers can cite them without copying bib blocks.
% These refer to the companion archive already-published.zip.

\bibitem{series:schedcompiler}
Scheduling and Compiler Techniques for Metadata-Hiding Overlay Lookups.
\newblock Anonymity series paper (published; companion archive \texttt{already-published.zip}), 2026.
\newblock Wiki: \texttt{[[2026.01.22 - Mathematics: Scheduling and Compiler Techniques for Metadata-Hiding Overlay Lookups]]}.

\bibitem{series:prefixcapacities}
Prefix Capacities and Separation Cuts for Shared-Kernel Termination-Time Hiding.
\newblock Anonymity series paper (published; companion archive \texttt{already-published.zip}), 2026.
\newblock Wiki: \texttt{[[2026.01.25 - Mathematics: Prefix Capacities and Separation Cuts for Shared-Kernel Termination-Time Hiding]]}.

\bibitem{series:spectral}
Spectral Anonymity and Optimal Distinguishing Bounds for Random-Walk Delegation in (Possibly Directed) Overlays.
\newblock Anonymity series paper (published; companion archive \texttt{already-published.zip}), 2026.
\newblock Wiki: \texttt{[[2026.01.23 - Mathematics: Spectral Anonymity and Optimal Distinguishing Bounds for Random-Walk Delegation in (Possibly Directed) Overlays]]}.

\bibitem{series:stoptimeaddendum}
Stop-Time Padding Addendum (unique-only): Max-Mixture Separations and Tail-Sign Deadline Normal Forms.
\newblock Anonymity series paper (published; companion archive \texttt{already-published.zip}), 2026.
\newblock Wiki: \texttt{[[2026.01.26 - Mathematics: Stop-Time Padding Addendum, Max-Mixture Separations and Tail-Sign Deadline Normal Forms]]}.

\bibitem{SeriesManyTesting}
Many-Testing Lower Bounds for Low-Latency Privacy on Hierarchical Overlays.
\newblock Anonymity series paper (published; companion archive \texttt{already-published.zip}), 2026.

\bibitem{series:knapsacktransport}
Optimal Poset-Feasible Padding: Knapsack, Transport, and Dual Certificates for TV Privacy.
\newblock Anonymity series paper (published; companion archive \texttt{already-published.zip}), 2026.
\newblock Wiki: \texttt{[[2026.01.29 - Mathematics: Optimal Poset-Feasible Padding: Knapsack, Transport, and Dual Certificates for TV Privacy]]}.
% Common bibliography items for draft papers in the anonymity series.
% These are used for cross-references between the split papers.

\bibitem{series:profiling}
Anonymous DHT Profiling and the Equalization Contract.
\newblock Anonymity series Paper~B, 2026. (This archive.)

\bibitem{series:abomoinl}
ABOM and OINL for Anonymity Claims: Proof-Carrying Odds-Inflation Budgets Without Publishing Full Artifacts.
\newblock Anonymity series Paper~C, 2026. (This archive.)

\bibitem{series:releaseA}
Privacy as a Release Property: Proof-Carrying Anonymity Receipts for Anonymous DHTs.
\newblock Release and Destination series Paper~A, 2026. (This archive.)


\bibitem{MironovRDP2017}
I.~Mironov.
\newblock R\'enyi differential privacy.
\newblock In \emph{IEEE Computer Security Foundations Symposium (CSF)}, 2017. arXiv:1702.07476.

\bibitem{DworkRoth2014}
C.~Dwork and A.~Roth.
\newblock The algorithmic foundations of differential privacy.
\newblock \emph{Foundations and Trends in Theoretical Computer Science}, 9(3--4), 2014.

\bibitem{KiferMachanavajjhala2014Pufferfish}
D.~Kifer and A.~Machanavajjhala.
\newblock A framework for mathematical privacy definitions.
\newblock \emph{ACM Transactions on Database Systems}, 39(1):3:1--3:36, Feb.\ 2014.
\newblock doi:10.1145/2514681. (PDF: \url{https://users.cs.duke.edu/~ashwin/pubs/pufferfish_TODS.pdf}).

\bibitem{LiuCuffVerdu2017}
J.~Liu, P.~Cuff, and S.~Verd\'u.
\newblock $E_\gamma$-resolvability.
\newblock \emph{IEEE Transactions on Information Theory}, 63(5):2629--2658, May 2017.
\newblock doi:10.1109/TIT.2016.2642111. Preprint: arXiv:1511.07829.

\bibitem{NuradhaMeasuredHockeyStick2025}
T.~Nuradha, V.~Singh, and M.~M.~Wilde.
\newblock Measured hockey-stick divergence and its applications to quantum pufferfish privacy.
\newblock In \emph{Proc. IEEE International Symposium on Information Theory (ISIT)}, 2025.
\newblock doi:10.1109/ISIT63088.2025.11195501. arXiv:2501.12359. doi:10.48550/arXiv.2501.12359.
\newblock \url{https://arxiv.org/abs/2501.12359}.

\bibitem{Weissman2003L1}
T.~Weissman, E.~Ordentlich, G.~Seroussi, S.~Verd\'u, and M.~J.~Weinberger.
\newblock Inequalities for the $L_1$ deviation of the empirical distribution.
\newblock Hewlett-Packard Labs Technical Report HPL-2003-97R1, 2003.
\newblock (PDF: \url{https://www.hpl.hp.com/techreports/2003/HPL-2003-97R1.pdf}; mirror: \url{https://shiftleft.com/mirrors/www.hpl.hp.com/techreports/2003/HPL-2003-97R1.pdf}).




\bibitem{series:state1}
Anonymous.
\newblock State-Dependent Anonymity in Anonymous DHTs.
\newblock AnonDHT state series Paper~1, 2026. (This archive.)

\bibitem{series:operationalA}
Anonymous.
\newblock Schedule-Only Retries in Anonymous DHTs: SRP and Retry-Trace Equalization.
\newblock Operational series Paper~A, 2026. (This archive.)

\bibitem{series:operationalB}
Anonymous.
\newblock Auditable Trace Privacy with Abort.
\newblock Operational series Paper~B, 2026. (This archive.)

\bibitem{series:eval1}
Anonymous.
\newblock Notarized Anonymity Certificates: Attempt Logs, Spending Discipline, and Verifier Rules.
\newblock Evaluation series Paper~1, 2026. (This archive.)

\bibitem{series:artifactpolicy}
Anonymous.
\newblock Artifact policy (source-only archive; artifacts available on request).
\newblock Series note, 2026.


\bibitem{series:notation}
Anonymous.
\newblock Notation and Minimal Model for the One-True-Archive.
\newblock Series synthesis note (Synthesis~8), 2026. (This archive.)

\bibitem{series:bossfightA}
Anonymous.
\newblock Boss Fight Budgets: Maximal Leakage as a Repeated-Use Anonymity Contract.
\newblock Boss Fight series Paper~A, 2026. (This archive.)

\bibitem{series:workedexample}
Anonymous.
\newblock Worked Example: Instantiating a Tiered Reader-Privacy Receipt.
\newblock Synthesis~17, The-One-True-Archive (2026).

\end{thebibliography}

\end{document}
